\section{VMA 权限标志与 PTE 标志的关系}
\label{sec:vma-flags}

\codefile{src/include/kernel/vma.h}
\begin{lstlisting}[style=cstyle]
#define VMA_READ    0x01u   /* 可读 */
#define VMA_WRITE   0x02u   /* 可写 */
#define VMA_EXEC    0x04u   /* 可执行（配合 NX 位使用） */
\end{lstlisting}

VMA 标志描述的是\emph{语义层级}的权限意图；PTE 标志是\emph{硬件层级}的访问控制。
它们之间的映射关系如表~\ref{tab:vma-pte-flags}~所示。

\begin{table}[H]
\centering
\begin{tabular}{lll}
\toprule
\textbf{VMA 标志} & \textbf{对应 PTE 操作} & \textbf{说明} \\
\midrule
\texttt{VMA\_READ}  & \texttt{PTE\_PRESENT}   & 页面存在即可读 \\
\texttt{VMA\_WRITE} & \texttt{PTE\_WRITABLE}  & 设置 R/W 位 \\
\texttt{VMA\_EXEC}  & 清除 NX 位（PAE/64-bit） & 32 位无 NX，留作扩展 \\
\bottomrule
\end{tabular}
\caption{VMA 权限标志到 PTE 标志的映射}
\label{tab:vma-pte-flags}
\end{table}

\begin{thinkbox}
在 32 位非 PAE 模式下，x86 没有 NX（No-Execute）位——所有可读页面都可以执行。
这意味着 \texttt{VMA\_EXEC} 在当前实现中其实不会产生硬件层面的约束。
那为什么还要定义它？

提示：考虑将来迁移到 PAE 或 64 位模式、以及软件层面的审计需求。
\end{thinkbox}

% ============================================================================

